\documentclass[tikz,border=10pt]{standalone}
\usepackage{tikz}
\usepackage{pgfplots}
\pgfplotsset{compat=1.18}
\usetikzlibrary{3d,calc}

\begin{document}

% ============================================
% 图1：偏导数的几何意义 - 3D曲面与切面
% ============================================
\begin{tikzpicture}[scale=1.2]
    \begin{axis}[
        width=12cm,
        height=10cm,
        view={120}{30},
        xlabel={$x$},
        ylabel={$y$},
        zlabel={$z$},
        xlabel style={sloped like x axis},
        ylabel style={sloped like y axis},
        zlabel style={sloped like z axis},
        xmin=-3, xmax=3,
        ymin=-3, ymax=3,
        zmin=0, zmax=18,
        colormap/viridis,
        colorbar,
        colorbar style={
            title={$z$值},
            yticklabel style={font=\small}
        },
        title={\textbf{偏导数的几何意义：} $z = x^2 + y^2$},
        title style={font=\large, yshift=0.5cm}
    ]
    
    % 绘制曲面 z = x^2 + y^2
    \addplot3[
        surf,
        shader=interp,
        opacity=0.7,
        domain=-2.5:2.5,
        domain y=-2.5:2.5,
        samples=30,
        samples y=30,
    ] {x^2 + y^2};
    
    % 标记点 P(1, 1, 2)
    \addplot3[only marks, mark=*, mark size=3pt, color=red] coordinates {(1, 1, 2)};
    \node[pin=120:{$P(1,1,2)$}] at (axis cs:1,1,2) {};
    
    % 沿x方向的切面（固定y=1）
    \addplot3[
        mesh,
        color=red,
        thick,
        domain=-2:2,
        domain y=1:1,
        samples=20,
    ] {x^2 + 1};
    
    % 沿y方向的切面（固定x=1）
    \addplot3[
        mesh,
        color=blue,
        thick,
        domain=1:1,
        domain y=-2:2,
        samples=20,
    ] {1 + y^2};
    
    % 添加图例说明
    \node[anchor=north west, fill=white, draw=black, rounded corners] 
        at (axis cs:2.5,2.5,15) {
        \begin{tabular}{l}
            \textcolor{red}{$\bullet$} 点 $P(1,1,2)$ \\
            \textcolor{red}{$-$} 沿$x$方向切面 ($y=1$) \\
            \textcolor{blue}{$-$} 沿$y$方向切面 ($x=1$) \\
        \end{tabular}
    };
    
    \end{axis}
\end{tikzpicture}

\newpage

% ============================================
% 图2：梯度下降过程可视化
% ============================================
\begin{tikzpicture}[scale=1.2]
    \begin{axis}[
        width=12cm,
        height=10cm,
        view={60}{30},
        xlabel={$w_1$},
        ylabel={$w_2$},
        zlabel={$L$},
        xlabel style={sloped like x axis},
        ylabel style={sloped like y axis},
        zlabel style={sloped like z axis},
        xmin=-3, xmax=3,
        ymin=-3, ymax=3,
        zmin=0, zmax=10,
        colormap/hot,
        title={\textbf{梯度下降：}在损失函数曲面上寻找最小值},
        title style={font=\large, yshift=0.5cm}
    ]
    
    % 绘制损失函数曲面（碗状）
    \addplot3[
        surf,
        shader=interp,
        opacity=0.6,
        domain=-2.5:2.5,
        domain y=-2.5:2.5,
        samples=25,
        samples y=25,
    ] {0.5*(x^2 + y^2)};
    
    % 梯度下降路径点
    \def\points{}
    \foreach \i in {0,1,...,10} {
        \pgfmathsetmacro{\t}{2.5*0.8^\i}
        \pgfmathsetmacro{\z}{0.5*(\t^2 + \t^2)}
        \addplot3[only marks, mark=*, mark size=2pt, color=red] 
            coordinates {(\t,\t,\z)};
    }
    
    % 最低点
    \addplot3[only marks, mark=*, mark size=4pt, color=green] 
        coordinates {(0,0,0)};
    \node[pin=-90:{\textcolor{green}{最小值点}}] at (axis cs:0,0,0) {};
    
    % 起始点标注
    \node[pin=45:{\textcolor{red}{起始点}}] at (axis cs:2.5,2.5,6.25) {};
    
    % 添加箭头表示梯度方向
    \draw[->, thick, red] (axis cs:2.5,2.5,6.25) -- (axis cs:2,2,4);
    \draw[->, thick, red] (axis cs:2,2,4) -- (axis cs:1.6,1.6,2.56);
    \draw[->, thick, red] (axis cs:1.6,1.6,2.56) -- (axis cs:1.28,1.28,1.638);
    
    \end{axis}
\end{tikzpicture}

\newpage

% ============================================
% 图3：链式法则的分解图
% ============================================
\begin{tikzpicture}[
    node distance=2.5cm,
    box/.style={rectangle, draw, fill=blue!20, thick, 
                minimum width=2cm, minimum height=1cm, font=\large},
    arrow/.style={->, thick, >=stealth},
    label/.style={font=\small, fill=white, inner sep=2pt}
]

% 标题
\node[font=\Large\bfseries] at (0, 3) {链式法则：复合函数求导的分解};

% 节点
\node[box] (x) at (-6, 0) {$x$};
\node[box] (y) at (-2, 0) {$y = f(x)$};
\node[box] (z) at (2, 0) {$z = g(y)$};
\node[box] (dzdx) at (6, 0) {$\frac{dz}{dx}$};

% 箭头
\draw[arrow] (x) -- (y) node[midway, above, label] {$\frac{dy}{dx}$};
\draw[arrow] (y) -- (z) node[midway, above, label] {$\frac{dz}{dy}$};
\draw[arrow] (z) -- (dzdx) node[midway, above, label] {相乘};

% 公式
\node[font=\Large] at (0, -2) {$\displaystyle \frac{dz}{dx} = \frac{dz}{dy} \cdot \frac{dy}{dx}$};

% 说明
\node[text width=10cm, align=center, font=\small] at (0, -3.5) {
    从 $x$ 到 $z$ 的变化率 = 从 $x$ 到 $y$ 的变化率 $\times$ 从 $y$ 到 $z$ 的变化率
};

\end{tikzpicture}

\newpage

% ============================================
% 图4：神经网络中的反向传播
% ============================================
\begin{tikzpicture}[
    node distance=1.8cm,
    neuron/.style={circle, draw, fill=blue!30, thick, minimum size=1cm, font=\large},
    input/.style={circle, draw, fill=green!30, thick, minimum size=1cm, font=\large},
    output/.style={circle, draw, fill=red!30, thick, minimum size=1cm, font=\large},
    arrow/.style={->, thick, >=stealth},
    gradarrow/.style={<-, thick, >=stealth, dashed, red},
    label/.style={font=\small}
]

% 标题
\node[font=\Large\bfseries] at (0, 4) {神经网络中的链式法则（反向传播）};

% 前向传播节点
\node[input] (x) at (-4, 0) {$x$};
\node[neuron] (z) at (-1, 0) {$z$};
\node[neuron] (a) at (2, 0) {$a$};
\node[output] (L) at (5, 0) {$L$};

% 前向箭头
\draw[arrow] (x) -- (z) node[midway, above, label] {$z = wx+b$};
\draw[arrow] (z) -- (a) node[midway, above, label] {$a = \sigma(z)$};
\draw[arrow] (a) -- (L) node[midway, above, label] {$L = (a-y)^2$};

% 反向传播箭头（梯度）
\draw[gradarrow] (x) -- (z) node[midway, below, label] {$\frac{\partial z}{\partial w}=x$};
\draw[gradarrow] (z) -- (a) node[midway, below, label] {$\frac{\partial a}{\partial z}$};
\draw[gradarrow] (a) -- (L) node[midway, below, label] {$\frac{\partial L}{\partial a}$};

% 总梯度公式
\node[font=\large] at (0, -3) {
    $\displaystyle \frac{\partial L}{\partial w} = \frac{\partial L}{\partial a} \cdot \frac{\partial a}{\partial z} \cdot \frac{\partial z}{\partial w}$
};

% 图例
\node[anchor=west] at (-6, -4) {\textcolor{blue}{$\bullet$} 前向传播};
\node[anchor=west] at (-6, -4.5) {\textcolor{red}{$\bullet$} 反向传播（梯度）};

\end{tikzpicture}

\newpage

% ============================================
% 图5：导数与微分的几何对比
% ============================================
\begin{tikzpicture}[scale=1.5]
    % 坐标轴
    \draw[->] (-0.5, 0) -- (5, 0) node[right] {$x$};
    \draw[->] (0, -0.5) -- (0, 4) node[above] {$y$};
    
    % 函数曲线 y = x^2 / 2
    \draw[thick, blue, domain=0.3:4.5, smooth, variable=\x] 
        plot (\x, {0.5*\x*\x/2});
    \node[blue] at (4.5, 2.5) {$y = f(x)$};
    
    % 点 P(2, 1)
    \fill[red] (2, 1) circle (2pt);
    \node[red, above right] at (2, 1) {$P(x_0, f(x_0))$};
    
    % 切线
    \draw[thick, red, dashed, domain=0.5:3.5] 
        plot (\x, {0.5*\x - 0});
    \node[red] at (3.5, 1.8) {切线};
    
    % dx 和 dy 标注
    \draw[<->, thick, green!60!black] (2, 1) -- (3, 1) node[midway, below] {$dx$};
    \draw[<->, thick, orange] (3, 1) -- (3, 1.5) node[midway, right] {$dy$};
    
    % Δx 和 Δy 标注
    \draw[<->, thick, green!60!black, dashed] (2, 1) -- (3, 1);
    \draw[<->, thick, purple, dashed] (3, 1) -- (3, {0.5*9/2}) node[midway, right] {$\Delta y$};
    \fill[purple] (3, {0.5*9/2}) circle (1.5pt);
    
    % 公式
    \node[font=\large] at (2.5, -1.2) {
        $\displaystyle \frac{dy}{dx} = f'(x_0) \quad \Rightarrow \quad dy = f'(x_0) \cdot dx$
    };
    
    \node[font=\small, text width=8cm, align=center] at (2.5, -2.2) {
        $\Delta y$ = 真实变化（曲线上）\\
        $dy$ = 线性近似（切线上）\\
        当 $dx \to 0$ 时：$\Delta y \approx dy$
    };
    
\end{tikzpicture}

\end{document}
